\subsection{The $N=2$ Supersymmetric Lagrangian for Gauge Fields}

The on-shell $N=1$ scalar multiplet $(A, \psi)$ and vector multiplet
$(A_\mu,\lambda)$, put together, have the same field content as the
on-shell $N=2$ vector multiplet $(A,\psi,\lambda, A_\mu)$. The
Lagrangian (\ref{fullC}) contains all these fields but as such is not
$N=2$ supersymmetric. We now assume that $(A,\psi,\lambda, A_\mu)$
form an $N=2$ vector multiplet and discuss the restrictions which this
assumtion imposes on the $N=1$ Lagrangian in (\ref{fullC}). First,
since $A^a_\mu$ and $\lambda^a$ belong to the adjoint representation
of the gauge group, $A_i$ and $\psi_i$ should also belong to the same
representation if they are to be part of the same multiplet. Hence,
$T^a_{ij}=-if^a_{ij}$ and the sets of indices $\{ i \}$ and $\{ a \}$
coincide. Second, since the two supersymmetry generators in the $N=2$
algebra appear on the same footing, the same must be the case with the
fermions $\psi^a$ and $\lambda^a$ in (\ref{fullC}). To satisfy this
condition, we set the superpotential ${\cal W}$ to zero since it
couples only to $\psi^a$. This condition also fixes the arbitrary
relative normalization between the Yang-Mills part and the scalar part
of the Lagrangian since it requires that the kinetic energy terms for
both fermions should have the same normalization.  This is achieved by
scaling $\Phi\rightarrow\Phi/g$ in (\ref{fullC}).  It turns out that
if the Lagrangian (\ref{fullC}) satisfies these conditions, then it
has $N=2$ supersymmetry. The terms containing the auxiliary fields now
take the form
$$
\frac{1}{g^2}{\rm Tr}\left(\frac{1}{2}DD +D\,[A^\dagger\,,\,A] + 
F^\dagger F \right)\,, 
$$
where, we have used the notation $\Phi=\Phi^a T^a$, with $T^a$ in the
fundamental representation. On eliminating $D$ and $F$, we get the
scalar potential    
\begin{equation}
V=- \frac{1}{2g^2} {\rm Tr}\left([A^\dagger\,,\,A]^2 \right)\,.
\label{scalarV}
\end{equation}
The full Lagrangian with $N=2$ supersymmetry can now be written as  
\begin{eqnarray}
{\cal L}&=&\frac{1}{8\pi}{\rm Im\, Tr}\left[\,\tau\left(\int d^2\theta\,
W^\alpha W_\alpha +2\,\int d^2\theta d^2\bar\theta\, \Phi^\dagger 
e^{-2V}\Phi \right)\right] 
\nonumber \\ 
&=&\frac{1}{g^2}{\rm Tr}\Big(-\frac{1}{4}F_{\mu\nu}F^{\mu\nu}+ 
g^2\frac{\theta}{32\pi^2}F_{\mu\nu}\widetilde F^{\mu\nu}+
(D_\mu A)^\dagger D^\mu A -\frac{1}{2} [A^\dagger\,,\,A]^2
\nonumber\\
&&{\hphantom {\frac{1}{g^2}{\rm Tr}(\,} }
-i\,\lambda\sigma^\mu D_\mu\bar\lambda -i\,\bar{\psi} 
\bar{\sigma}^\mu D_\mu\psi-i\sqrt{2}\,[\lambda,\psi]\,A^\dagger
-i\sqrt{2}\,[\bar\lambda,\bar\psi]\,A \Big)\,,
\label{N=2}
\end{eqnarray}
where, in the component expansion, the auxiliary fields have been
eliminated. 
Note that (\ref{N=2}) is the supersymmetric generalization of the
Yang-Mills-Higgs theories described in section 1, with the Higgs
potential $V$ given by (\ref{scalarV}), and with a $\theta$-term.    
Therefore, the discussion in section 1 also applies to our $N=2$
supersymmetric theory. In particular, the Higgs vacuum is defined by
$D_\mu A=0, V=0$, and the potential vanishes for non-zero field
configurations provided $A$ commutes with $A^\dagger$. Thus the model
admits monopole and dyonic solutions and 
contains massive gauge bosons. For example, if the gauge group is
$SU(2)$ or $SO(3)$, it is broken down to $U(1)$ and two of the gauge
bosons become massive. 

Now, suppose that we are interested in the behaviour of this model at
energies lower than some cutoff $\Lambda$ which is smaller than the
mass of the lightest massive state in the theory. At such energies, we
will not encounter any on-shell massive states and the physics can be
described by the Wilsonian low-energy effective action. This effective
action is obtained by completely integrating out all massive states as
well as integrating out all massless excitations above the scale
$\Lambda$ \cite{SV1,SV2}. This, in general, is a complicated procedure and
cannot be carried out explicitly. Fortunately, in our model, the
general form of the Wilsonian effective action is severely constrained
by $N=2$ supersymmetry. This is easiest to see when the theory
(\ref{N=2}) is formulated in term of $N=2$ superfields as described
below. 
 
\noindent\underline{The $N=2$ Superspace Formulation:}
The $N=2$ superspace is obtained by adding four more fermionic degrees
of freedom, say, $\tilde\theta$ and $\bar{\tilde\theta}$, 
to the $N=1$ superspace. Thus, a generic $N=2$ superfield can be
written as $F(x,\theta,\bar\theta,\tilde\theta,\bar{\tilde\theta})$. 
We need a superfield which has the same components as the $N=2$ vector
multiplet. This is obtained by imposing the constraints of chirality
and reality on a general $N=2$ superfield \cite{Grimm}. We briefly
describe this following the approach of \cite{bilal}. An $N=2$ chiral
superfield $\Psi$ is defined by the constraints $\bar
D_{\dot\alpha}\Psi=0$ and $\bar{\tilde D}_{\dot\alpha}\Psi=0$, where,
the supercovariant derivative $\bar{\tilde D}_{\dot\alpha}$ is defined
in the same way as $\bar D_{\dot\alpha}$ with $\theta$ replaced by
$\tilde\theta$. The expansion of $\Psi$ can be arranged in powers of
$\tilde\theta$ and can be written as
$$
\Psi=\Psi^{(1)}(\tilde y,\theta)+\sqrt{2}\tilde\theta^\alpha
\Psi^{(2)}_\alpha(\tilde y, \theta) + \tilde\theta^\alpha
\tilde\theta_\alpha \Psi^{(3)}(\tilde y, \theta)\,,
$$
where, $y^\mu =x^\mu +i\theta\sigma^\mu\bar\theta +i\tilde\theta
\sigma^\mu\bar{\tilde\theta}$. This expansion helps us relate the
$N=2$ formalism to the $N=1$ language we have been using so far.
Clearly, the component $\Psi^{(1)}$ has the same form as the $N=1$
chiral superfield $\Phi$. The remaining two components are constrained
by a reality condition. The outcome is that $\Psi^{(2)}_\alpha=
W_\alpha (\tilde y, \theta)$ as given in (\ref{W-nonabelian}), and
$\Psi^{(3)}$ is given by
$$
\Psi^{(3)}(\tilde y,\theta)= \Phi^\dagger(\tilde y-i\theta\sigma 
\bar\theta, \theta,\bar\theta) exp \left[2gV(\tilde y -i\theta\sigma
\bar\theta, \theta,\bar\theta\right]\,\mid_{\bar\theta\bar\theta}\,.
$$
Here, $\Phi(\tilde y -i\theta\sigma\bar\theta,\theta,\bar\theta)$ is
to be understood as the expansion in (\ref{csfield}). Clearly, $\Psi$
has the same field content as the $N=2$ vector multiplet. One can
verify that in terms of the $N=2$ superfield $\Psi$, the $N=2$
Lagrangian (\ref{N=2}) is given by the compact expression
\begin{equation}
{\cal L}= \frac{1}{4\pi}{\rm Im\,Tr}\int d^2\theta d^2\tilde\theta\,
\frac{1}{2}\tau\Psi^2\,.
\label{N=2manifest}
\end{equation}

Using the $N=2$ chiral superfield $\Psi$ (also referred to as the
$N=2$ vector superfield), we can now construct the most general
$N=2$ Lagrangian for the gauge fields: Corresponding to any function
${\cal F}(\Psi)$, we can construct a Lagrangian  
\begin{eqnarray}
{\cal L}&=& \frac{1}{4\pi}{\rm Im\,Tr}\int d^2\theta d^2\tilde\theta
\,{\cal F}(\Psi) 
\nonumber\\
&=&\frac{1}{8\pi}{\rm Im}\left(\int d^2\theta\, {\cal F}_{ab}(\Phi) 
W^{a\alpha}W^b_\alpha + 2\int d^2\theta d^2\bar\theta\, (\Phi^\dagger 
e^{2gV})^a {\cal F}_a(\Phi)\right)\,. 
\label{N=2general}
\end{eqnarray}
Here, ${\cal F}_a(\Phi)=\partial{\cal F}/\partial\Phi^a$,
${\cal F}_{ab}(\Phi)=\partial^2{\cal F}/\partial\Phi^a\partial\Phi^b$
and ${\cal F}$ is referred to as the $N=2$ prepotential. From the
above, we can easily read off the K\"{a}hler potential as
${\rm Im}(\Phi^{\dagger a} {\cal F}_a(\Phi)$. This gives rise to a metric 
$g_{ab}={\rm Im}(\partial_a\partial_b{\cal F})$ on the space of fields. 
A metric of this form is called a special K\"{a}hler metric.
If we demand renormalisability, then ${\cal F}$ has to be quadratic in
$\Psi$ as in (\ref{N=2manifest}). However, if we want to write a
low-energy effective action, then renormalisability is not a criterion
and ${\cal F}$ can have a more complicated form. In particular, we can
start from the microscopic theory (\ref{N=2}), corresponding to a
quadratic prepotential, and try to construct the modified ${\cal F}$
for the low-energy Wilsonian effective action. The exact determination
of this function is the subject of the work of Seiberg and Witten.

\subsection{The $N=2$ Supersymmetric Lagrangian for Matter Fields}

Since matter fields and gauge fields transform under different
representations of the gauge group, they cannot be part of the same
multiplet. The $N=2$ matter supermultiplet is called the
hypermultiplet and contains one pair of complex scalars and one pair
of two-component spinors, all transforming under the same
representation of the gauge group. From our discussion of the
representations of the supersymmetry algebra it follows that if a
hypermultiplet is massive, then its mass should appear as a central
extension in the supersymmetry algera. This is so because otherwise
$N=2$ supersymmetry requires a larger number of components than is
contained in a hypermultiplet. In $N=1$ notation, a hypermultiplet
contains a chiral superfield $Q$ and an anti-chiral superfield
$\widetilde Q^\dagger$, both transforming under the same representatin
$N_c$ of the gauge group $SU(N_c)$ We denote the components of $Q$ and
$\widetilde Q$ by $(q,\psi_q,F_q)$ and $(\widetilde q,
\psi_{\widetilde q},F_{\widetilde q})$ respectively. The form of the
Lagrangian for $N_f$ hypermultiplets (labelled by an index $i$),
interacting with a $N=2$ vector multiplet, can be partly inferred from
the $N=1$ theories and is given by
\begin{eqnarray}
{\cal L}=&&\int d\theta^4\left( Q^\dagger_i e^{-2V}Q_i +\widetilde Q_i
e^{2V}\widetilde Q^\dagger_i\right)+\int
d\theta^2\left(\sqrt{2}\widetilde Q_i\Phi Q_i+m_i\widetilde Q_i
Q_i\right)+\,h.c.\nonumber\\       &&+\cdots
\label{N=2hyper}
\end{eqnarray}
Here, the dots represent the Lagrangian for the pure $N=2$ vector
multiplet and we have supressed the gauge group indices. The term  
$\sqrt{2}\widetilde Q_i\Phi Q_i$ is related, by $N=2$ supersymmetry, to
the coupling of the hypermultiplet with the $N=1$ vector multiplet
$V$. The presence of this term is not required in an $N=1$ Lagrangian.  
As in normal QCD, when all masses $m_i$ are equal, the theory is
invariant under the global flavour group $SU(N_f)$. 

Eliminating the auxiliary fields $F_q$ and $F_{\widetilde q}$, which
appear in the hypermultiplet, results in a contribution to the scalar
potential given by (compare with (\ref{Spot}))
$$
V=\frac{1}{2}g^2\sum_a D_a D^a\,,
$$
with
\begin{equation}
D^a=\sum_{i=1}^{N_f}\left(q^\dagger_i\lambda^a q_i -\widetilde
q_i\lambda^a \widetilde q_i^\dagger\right)\,.
\label{D-term}
\end{equation}
Here, $\lambda^a$ are the gauge group generators in the fundamental
representaion. This term is referred to as the $D$-term. It must be
mentioned that the $N=2$ algebra has a global $SU(2)_R$ symmetry which
should also be a symmetry of the Lagrangian. However, in our
decomposition of the hypermultiplet in terms of $Q$ and $\widetilde
Q^\dagger$, this symmetry is not manifest since it rotates the scalar
components $q$ and $\widetilde q^\dagger$ as a doublet. It is not
difficult to write an $N=2$ Lagrangian with manifest $SU(2)_R$
symmetry \cite{West}. 
 
\noindent\underline{Minimization of the $D$-term}:
To find the vacuum of the theory, now we also have to minimize the
$D$-term contribution and the hypermultiplet mass term contribution to
the scalar potential \cite{DDS,ADS}. For non-zero quark masses, the
only solution is 
$q=\widetilde q =0$. Thus, only the scalar $A$ can have a non-zero
vacuum expectation value. But when $m_i=0$, then the $D$-term can be
minimized for a set of non-vanishing $q$ and $\widetilde q$ (the
potential has flat directions). However, the $\widetilde Q\Phi Q$ term
in the Lagrangian now requires $A_{vac}=0$. In the following, we
determine these flat directions.

The minimum of the $D$-term corresponds to $D^a=0$. To solve this,
note that $q^{(i)}_\alpha$ (where $\alpha$ denotes a colour index) can
be regarded as a set of $N_c$ vectors in ${\bf C}^{N_f}$. Thus, we can
construct the matrix of scalar products, 
$$
\sum_i q^{(i)}_\alpha q^{\dagger (i)}_\beta=q_\alpha\cdot
q^\dagger_\beta=(qq^\dagger)_{\alpha\beta}\,.
$$
In terms of this,
$$
\sum_i q^\dagger_i\lambda^aq_i={\rm Tr}\,\left(qq^\dagger
\lambda^a\right) \,.
$$
Similarly, we construct the matrix $\widetilde q_\alpha^\dagger \cdot
\widetilde q_\beta = (\widetilde q^\dagger\widetilde q)_{\alpha\beta}$.
Hence, $D^a=0$ can be written as 
$$
D^a={\rm Tr}\left[\left(qq^\dagger-\widetilde q^\dagger
\widetilde q\right)\,\lambda^a\right]=0 \,.
$$
Since $\lambda^a$ are in an irreducible representation of $SU(N_c)$,
the above condition is solved by 
\begin{equation}
(qq^\dagger) - (\widetilde q^\dagger\widetilde q)=c^2{\bf 1}_{N_c}\,.
\label{D=0} 
\end{equation}
Here, we can distinguish two cases: $N_f < N_c$ and $N_f\geq N_c$.

\noindent 1) $N_f < N_c$: In this case, the matrix $(qq^\dagger)$ has
rank $N_f$ and thus it has $N-N_f$ zero eigenvalues. The same applies to
the matrix $(\widetilde q^\dagger\widetilde q)$. By using $SU(N_c)\times
SU(N_f)\times U(1)_R$ rotations, $(qq^\dagger)$ can be diagonalized and
condition (\ref{D=0}) then implies that $(\widetilde
q^\dagger\widetilde q)$ must also be diagonal. Hence, in this basis, we
can solve for $q$ and $\widetilde q$ as
$$
q = \left(
\begin{array}{cccc}
v^{(1)}_1 & 0 &\cdots&0\\
0 & v^{(2)}_2 &&\\
\vdots&&\ddots & \\
&&&v^{(N_f)}_{N_f}\\
0&0&\cdots&0\\
\vdots&\vdots&&\vdots
\end{array}
\right)\,, \quad
\widetilde q = \left(
\begin{array}{cccccc}
\tilde v^{(1)}_1 & 0 & \cdots & & 0 & \cdots \\
 0 & \tilde v^{(2)}_2 & & & 0 & \cdots \\
\vdots  & & \ddots &  &   &  \\
 0 & & &\tilde v^{(N_f)}_{N_f} & 0 & \cdots  \\
\end{array}
\right)\,.
$$
The presence of zero eigenvalues imply that $c = 0$ and therefore,
$v^{(i)}_i =\tilde{v}^{(i)}_i$. In this case, the gauge symmetry is
broken down to $SU(N_c-N_f)$ except for $N_c=N_f-1$, where it is
totally broken. $2N_fN_c-N^2_f$ quark superfields become heavy and the
remaining $N^2_f$ quark superfields remain massless and correspond to
the Goldstone bosons of broken global symmetries.  

\noindent 2) $N_f \leq N$:  In this case $(qq^\dagger)$
has rank $N_c$ and, generically, its eigenvalues are not zero. As a
result, $c\not=0$. Arguing as above, in this case the vacuum values
of $q$ and $\widetilde q$ can be written as  
$$
q = \left(
\begin{array}{cccccc}
v^{(1)}_1 & 0 & \cdots & & 0 & \cdots \\
 0 & v^{(2)}_2 & & & 0 & \cdots \\
\vdots  & & \ddots &  &   &  \\
 0 & & & v^{(N_c)}_{N_c} & 0 & \cdots  \\
\end{array}
\right)\,,\quad
\widetilde q = \left(
\begin{array}{cccc}
\tilde v^{(1)}_1 & 0 &\cdots&0\\
0 &\tilde v^{(2)}_2 &&\\
\vdots&&\ddots & \\
&&&\tilde v^{(N_c)}_{N_c}\\
0&0&\cdots&0\\
\vdots&\vdots&&\vdots
\end{array}
\right)\,.
$$
and
$$
\tilde{v}_{i}^{(i)} = \sqrt{|v_{i}^{(i)}|^2 - c^2} \, .
$$
In this case, the gauge group is completely broken, and depending
on the values of $v_{i}^{(i)}$, the pattern of chiral symmetry
breaking quite complicated with many possibilities.
For $N_f > N$, there are surviving $R$-symmetries. 

\subsection{Central Charges in the $N=2$ Pure Gauge Theory}

We have seen that the $N=2$ supersymmetry algebra with central charge
$Z$ implies the bound $M\ge \sqrt{2}|Z|$ on the particle masses. It
was also stated that this bound is the same as the BPS bound on the
masses which is determined in terms of the electric and magnetic
charges. In this section, we prove the above statement by explicitly
calculating $Z$ following \cite{Olive-Witten}.  

In the supersymmetry algebra, The central charge $Z$ appears in the
commutator of the supercharges $Q^I_\alpha$ which, in turn, are space
integrals of $S^{I0}_\alpha$ (Here, $S^{I\mu}_\alpha$ denotes the
supercurrent). Thus, we first have to compute $S^{I0}$'s in terms of
the basic fields, and then evaluate their commutators. The central
charge is related to a surface term in the space integral of this
commutator which is non-zero if the field configuration corresponds to
electric and magnetic charges.
  
As a warm up exercise, we start with the theory for $N=1$ chiral
superfields. The Lagrangian is given by 
$$
{\cal L}=\int d^4 \theta\,\Phi^\dagger\Phi +\int d^2 \theta \;
{\cal W}(\Phi) + \int d^2 \bar\theta \bar{\cal W}(\Phi^\dagger)\,.
$$
Defining $y^\mu=x^\mu + i \theta \sigma^\mu \bar{\theta}$, the
superfield can be expanded as $\Phi=A(y)+\sqrt{2}\,\theta\psi(y)
+\theta\theta F(y)$. In this basis, $Q_{\alpha}=\partial /
\partial\theta^{\alpha}$, $\bar{Q}_{\dot{\alpha}}=-
\partial/\partial\bar{\theta}^{\dot{\alpha}} + 2i
(\theta \sigma^\mu)_{\dot{\alpha}}\,\partial/\partial y^\mu$, and the
supersymmetry variations, $\delta_{\epsilon}=\epsilon^\alpha Q_\alpha +
\bar{\epsilon}_{\dot{\alpha}} \bar{Q}^{\dot{\alpha}}$, of the fields
are given by 
$$
\begin{array}{ll}
\delta A=\sqrt{2} \epsilon \psi\,,&
\delta\bar A=\sqrt{2} \bar{\epsilon}\bar{\psi}\,, \\
\delta\psi=\sqrt{2}\epsilon F + i\sqrt{2}\sigma^\mu
\bar{\epsilon} \partial_\mu A\,, &
\delta\bar{\psi}=-i\sqrt{2}\epsilon\sigma^\mu
\partial_\mu \bar{A}+\sqrt{2} \bar{F} \bar{\epsilon}\,, \\
\delta F=i\sqrt{2} \bar{\epsilon}\bar{\sigma^\mu}
\partial_\mu \psi\,, & \delta \bar{F}=i\sqrt{2}\epsilon 
\sigma^\mu \partial_\mu \bar{\psi}\,.
\end{array}
$$
Using these transformations, we can compute the variation of the terms
involving the superpotential ${\cal W}$ as  
\begin{eqnarray}
\delta ({\cal W}'F - \frac{1}{2} \frac{\partial^2 {\cal W}}{\partial
  A^2}\psi\psi) 
&=& {\cal W}''\delta AF+{\cal W}'\delta F-\frac{1}{2}{\cal W}'''\delta
A\psi\psi+{\cal W}'' \psi \delta \psi \nonumber\\
&=& {\cal W}''\sqrt{2}\epsilon\psi F+i\sqrt{2}{\cal W}'\bar{\epsilon}
\bar{\sigma}^\mu \partial_\mu \psi+\sqrt{2}{\cal W}''\psi(i\sigma^\mu 
\bar{\epsilon}\partial_\mu A + \epsilon F) \nonumber\\
&=&\partial_\mu \left( i \sqrt{2} \frac{\partial {\cal W}}{\partial A}
\bar{\epsilon} \bar{\sigma}^\mu \psi \right )\,.\nonumber
\end{eqnarray}
Note that we have used ${\cal W}'''\delta A\psi\psi={\cal W}'''
\sqrt{2}(\epsilon \psi)\psi\psi=0$. Since ${\cal W}'''$ is totally
symmetric, this statement is also true in the presence of many fields 
$\psi^i$. Thus, we have, 
$$
\delta (\int d^{2} \theta {\cal W} + \int d^{2} \bar{\theta}
\bar{{\cal W}})=i\sqrt{2}\partial_\mu\left(\frac{\partial {\cal W}}
{\partial A_j} \bar{\epsilon}\bar{\sigma}^\mu \psi^{j} -
\frac{\partial \bar{{\cal W}}}{\partial \bar{A}_j} \bar{\psi}_{j}
\bar{\sigma}^\mu \epsilon \right)\,.
$$
Next we calculate the variation of the kinetic terms,
$$
{\cal L}_{D} = \partial_\mu A^\dagger \partial^\mu A +
F^\dagger F-\frac{i}{2}\bar{\psi} \bar{\sigma}^\mu\partial_\mu\psi+
\frac{i}{2}\partial_\mu\bar{\psi}\bar{\sigma}^\mu\psi\,.
$$
A slightly lengthy computation yields:
\begin{eqnarray}
\delta {\cal L}_{D}=&-&\frac{i}{\sqrt{2}}\partial_\mu (F\bar{\psi}
\bar{\sigma}^\mu \epsilon ) +
\frac{i}{\sqrt{2}} \partial_\mu (\epsilon\psi\partial^\mu A^\dagger - 
\epsilon \sigma^{\nu\mu}\psi \partial_\nu A^\dagger)\nonumber \\
&+& \frac{i}{\sqrt{2}}\partial_\mu (F^\dagger \bar{\epsilon}
\bar{\sigma}^\mu\psi)+\frac{i}{\sqrt{2}}\partial_\mu (\bar{\epsilon} 
\bar{\psi}\partial^\mu A-\bar{\psi}\bar{\sigma}^{\mu\nu}\bar{\epsilon}
\partial_\nu A )\,.\nonumber
\end{eqnarray}
Adding the variations of the superpotential and the kinetic terms, and
using the definitions 
$$
\sigma^{\mu\nu}=\frac{1}{2} (\sigma^\mu \bar{\sigma}^\nu -
\sigma^\nu \bar{\sigma}^\mu )\,,\quad
\bar{\sigma}^{\mu\nu} = \frac{1}{2} (\bar{\sigma}^\mu \sigma^\nu
- \bar{\sigma}^\nu \sigma^\mu ) \,,
$$
we get the supercurrent as 
\begin{equation}
S^\rho_{matter}=\sqrt{2}\epsilon \sigma^\nu\bar{\sigma}^\rho \psi
\partial_\nu A^\dagger + i \sqrt{2}\frac{\partial\bar{{\cal W}}}
{\partial A^\dagger}\bar{\psi}\bar{\sigma}^\rho\epsilon\nonumber 
+ \sqrt{2}\bar{\psi} \bar{\sigma}^\rho\sigma^\nu\bar{\epsilon}
\partial_\nu A - i\sqrt{2}\frac{\partial {\cal W}}{\partial A} 
\bar{\epsilon} \bar{\sigma}^\rho \psi\,. 
\label{N=1scch}
\end{equation}
Note that for convenience, we have included the supersymmetry
transformation parameters $\epsilon,\bar\epsilon$ in the definition of
the supercurrent. So, expanding in components, we actually have:
\begin{equation}
S^\rho =\epsilon^\alpha S^\rho_\alpha + \bar\epsilon_{\dot\alpha}
\bar S^{\rho\dot\alpha}\,.
\label{S-comp}
\end{equation}
 
Next we consider the inclusion of gauge fields which are described by
a vector multiplet. The supersymmetry variations of the fields in this
multiplet are given by 
\begin{equation}
\begin{array}{lll}
\delta A^a_\mu &=& -i\bar{\epsilon}\bar{\sigma}_\mu\lambda^a +
i \bar{\lambda}^a \bar{\sigma}_\mu\epsilon \,,\\
\delta D^{a} &=&\bar{\epsilon}\bar{\sigma}^\mu D_\mu \lambda^a +
D_\mu \lambda^{a} \bar{\sigma}^\mu \epsilon \,,\\
\delta\lambda^a &=& \frac{1}{2} \sigma^{\mu\nu}\epsilon F^a_{\mu\nu} 
+ i\epsilon D^a \,,\\
\delta\bar{\lambda}^a &=&\frac{1}{2}\bar{\epsilon}\bar{\sigma}^{\nu
\mu} F^a_{\mu\nu} - i \bar{\epsilon} D^a \,.
\end{array}
\label{str1-g}
\end{equation}
Furthermore, in the presence of the gauge interactions, the matter
field transformations need some modification: 
\begin{equation}
\begin{array}{lll}
\delta A &=& \sqrt{2} \epsilon \psi \,,\\
\delta \psi &=& \sqrt{2} \epsilon F + i \sqrt{2} \sigma^\mu
\bar{\epsilon} D_\mu A \,,\\
\delta\bar{\psi}&=&\sqrt{2}\bar{\epsilon} F^\dagger -i\sqrt{2}
\epsilon \sigma^\mu D_\mu A^\dagger \,,\\
\delta F &=& i\sqrt{2}\bar{\epsilon} \bar{\sigma}^\mu D_\mu \psi - 
2i T^a A \bar{\epsilon} \bar{\lambda}^a\,.
\end{array}
\label{str1-m}
\end{equation}
The last term in $\delta F$ is needed to cancel part of the variation
of $\psi$ in the term $A^\dagger T^a\lambda^a\psi$ in the Lagrangian
(\ref{fullC}). The part of the Lagrangian describing the pure vector
superfieled can be written as 
$$
{\cal L}=\frac{1}{g^2}\left(-\frac{1}{4} F^a_{\mu\nu} F^{a\mu\nu}-
\frac{i}{2}\bar\lambda^a\bar\sigma^\mu D_\mu\lambda^a+\frac{i}{2}
D_\mu \bar\lambda^a\bar\sigma^\mu\lambda^a+\frac{1}{2}D^{2} \right) + 
\frac{\theta}{32 \pi^2} F^a_{\mu\nu} \tilde{F}^{a\mu\nu} \,.
$$
If we ignore the $\theta F \tilde{F}$ term, then the supercurrent
obtained from this part of the Lagrangian takes the form
$$
S^\rho_{gauge}=-\frac{i}{2g^2}\left(\bar{\lambda}^a \bar{\sigma}^\rho
\sigma^{\mu\nu}\epsilon F^a_{\mu\nu}+\bar{\epsilon}\bar{\sigma}^{\mu\nu} 
\bar{\sigma}^\rho \lambda^a F^a_{\mu\nu}\right)\,.
$$
For the theory coupled to matter, the supercurrent is not just the sum
of the above, plus the gauge-covariantized version of $S^\rho_{matter}$. 
We also expect some contribution from the terms $D^a A^\dagger T^a A$
and $i A^\dagger T^a \lambda^a\psi+h.c.$. Taking this into account,
the supercurrent for the interacting $N=1$ Lagrangian (\ref{fullC}) is
given by 
\begin{eqnarray}
S^\rho = &-&\frac{i}{2g^2}\left(\bar{\lambda}^a \bar{\sigma}^\rho
\sigma^{\mu\nu}\epsilon +\bar{\epsilon}\bar{\sigma}^{\mu\nu} 
\bar{\sigma}^\rho \lambda^a\right) F^a_{\mu\nu} -
\left(\bar{\epsilon} \bar{\sigma}^\rho \lambda^a + \bar{\lambda}^a 
\bar{\sigma}^\rho \epsilon\right) A^\dagger T^a A
\nonumber\\
&+& \sqrt{2}\epsilon\sigma^\mu \bar{\sigma}^\rho \psi D_\mu A^\dagger
+\sqrt{2}\bar{\psi}\bar{\sigma}^\rho\sigma^\mu\bar{\epsilon}D_\mu A
+i\sqrt{2}\frac{\partial \bar{\cal W}}{\partial A^\dagger}\bar{\psi}
\bar{\sigma}^\rho\epsilon-i\sqrt{2}\frac{\partial{\cal W}}{\partial A}
\bar{\epsilon}\bar{\sigma}^\rho \psi\,.
\label{N=1sc}
\end{eqnarray}

Now, let us turn our attention to the pure $N=2$ gauge theory. As
discussed in section 1, this is obtained from the $N=1$ theory
(\ref{fullC}) by setting ${\cal W}=0$ and by scaling the chiral superfield
$\Phi$ to $\Phi/g$. Furthermore, $\Phi$ is now a vector in the adjoint
representation of the gauge group. Thus, relabelling $(\lambda,\psi)$
as $(\lambda_1,\lambda_2)$, we can easily write down one of the $N=2$ 
supercurrents as  
\begin{eqnarray}
g^2 S^\rho_{(1)}=&-&\frac{i}{2}\left(\bar{\lambda}^a_1\bar{\sigma}^\rho
\sigma^{\mu\nu}\epsilon + \bar{\epsilon}\bar{\sigma}^{\mu\nu}
\bar{\sigma}^\rho \lambda^a_1\right) F^a_{\mu\nu} -
\left(\bar{\epsilon}\bar{\sigma}^\rho\lambda^a_{1}+\bar{\lambda}^a_{1}
\bar{\sigma}^\rho \epsilon\right) A^\dagger T^a A 
\nonumber \\
&+&\sqrt{2}\,\epsilon\sigma^\mu\bar{\sigma}^\rho\lambda^a_{2}\, D_\mu
A^{a\dagger}+\sqrt{2}\,\bar{\lambda}^a_{2}\bar{\sigma}^\rho \sigma^\mu
\bar{\epsilon}\, D_\mu A^a \,.\nonumber 
\end{eqnarray}

The $N=2$ theory is also invariant under a second set of supersymmetry
transformations (with parameter $\epsilon'$) which is obtained from
(\ref{str1-g}) and (\ref{str1-m}) by the replacement
$\lambda\rightarrow\psi$, $\psi\rightarrow -\lambda$. This corresponds
to a transformation of the ($\lambda,\psi$) doublet by an element of
$SU(2)_R$, which is a symmetry group of the $N=2$ algebra.  The
associated conserved current is then obtained from $S^\rho_{(1)}$ by
the replacement $\lambda_1 \rightarrow\lambda_2$,
$\lambda_2\rightarrow-\lambda_1$:
\begin{eqnarray}
g^2 S^\rho_{(2)}=&-&\frac{i}{2}\left(\bar{\lambda}^a_2\bar{\sigma}^\rho
\sigma^{\mu\nu}\epsilon' + \bar{\epsilon'}\bar{\sigma}^{\mu\nu}
\bar{\sigma}^\rho \lambda^a_2 \right) F^a_{\mu\nu} -
\left(\bar{\epsilon'}\bar{\sigma}^\rho\lambda^a_{2}+\bar{\lambda}^a_{2}
\bar{\sigma}^\rho \epsilon'\right) A^\dagger T^a A 
\nonumber \\
&-&\sqrt{2}\,\epsilon'\sigma^\mu \bar{\sigma}^\rho\lambda^a_{1}\,D_\mu
A^{a\dagger}-\sqrt{2}\,\bar{\lambda}^a_{1}\bar{\sigma}^\rho \sigma^\mu
\bar{\epsilon'}\, D_\mu A^a \,.\nonumber 
\end{eqnarray}
Let us first concentrate on $S^\mu_{(1)}$. Using the identities 
\begin{eqnarray}
\sigma^a \bar{\sigma}^b \sigma^c &=& \eta^{ab} \sigma^c - \eta^{ac}
\sigma^b + \eta^{bc} \sigma^a + i \epsilon^{abcd} \sigma_d \,, 
\nonumber\\
\bar{\sigma}^a \sigma^b \bar{\sigma}^c &=& \eta^{ab} \bar{\sigma}^c
- \eta^{ac} \bar{\sigma}^b + \eta^{bc} \bar{\sigma}^a - i
\epsilon^{abcd} \bar{\sigma}_d \nonumber\,,
\end{eqnarray}
along with $\chi\sigma^\mu\bar\psi=-\bar\psi\bar\sigma^\mu\chi$, this
supercurrent can be rewritten as 
\begin{eqnarray}
g^2S^\mu_{(1)}=&-&\epsilon\sigma_\nu\bar{\lambda}^a_{1}(iF^{a\mu\nu}+ 
\widetilde F^{a\mu\nu})+\sqrt{2}\,\epsilon\sigma^\nu\bar{\sigma}^\mu
\lambda^a_{2}\,D_\nu A^{\dagger a}+\epsilon\sigma^\mu 
\bar{\lambda}^a_{1} A^\dagger T^a A \nonumber \\
&+& (\mbox{\it $\bar\epsilon$ dependent terms})\,.\nonumber
\end{eqnarray}
From the above, we can easily read off the components
$S^\mu_{(1)\alpha}$ (see eq. (\ref{S-comp})) as
$$
g^2 S^\mu_{(1)\alpha}= \sigma_{\nu\alpha\dot\alpha}
\bar{\lambda}^{a\dot\alpha}_{1} (iF^{a\mu\nu}+\widetilde F^{a\mu\nu})
+\sqrt{2}(\sigma^\nu\bar{\sigma}^\mu\lambda^a_{2})_\alpha D_\nu 
A^{\dagger a}+\sigma^\mu_{\alpha\dot{\alpha}}
\bar{\lambda}^{a\dot{\alpha}}_{1} A^\dagger T^a A\,. 
$$
After lowering the spinor index using   
$$
\bar{\lambda}^{\dot{\alpha}}_{1}=\epsilon^{\dot{\alpha}\dot{\beta}}
\bar\lambda_{1\dot\beta}=i(\sigma_y\lambda^\dagger_1)^{\dot\alpha}\,,
\qquad(\mbox{Note that}\;\bar\lambda_{\dot\beta}=\lambda^\dagger_\beta) 
$$
and using vector notation for spatial components of vectors, 
The $\mu=0$ component of the current takes the form 
\begin{eqnarray}
g^2S^0_{(1)\alpha}\,\,=&-&i(\vec{\sigma}\sigma_y\lambda^{\dagger a}_1)_\alpha
\cdot (i\vec{F}^a+\vec{\widetilde F})^a \nonumber\\
&+& \sqrt{2} \lambda^a_{2\alpha} D_0 A^{\dagger a} +
\sqrt{2}(\vec{\sigma}\cdot\vec{D}A^{\dagger a}\lambda^a_{2})_\alpha 
+i(\sigma_y\lambda^{\dagger a}_1)_\alpha A^\dagger T^a A\,, \nonumber
\end{eqnarray}
where,
$$
\vec{F}^a=F^{a 0i}\,,\quad \vec{\widetilde F}^a=\widetilde F^{a0i}\,.
$$
As before, the expression for $S^0_{(2)\alpha}$ is obtained from the
above by the replacements $\lambda_1\rightarrow\lambda_2$, $\lambda_2
\rightarrow-\lambda_1$.

To evaluate the central charge, we are interested in the anti-commutator 
$$
\{Q_{(1)\alpha}\,,\,Q_{(2)\beta}\} =
\left\{\int d^3x\,S^0_{(1)\alpha}(\vec{x},0)\,,\,\int d^3y\,S^0_{(2)
\beta}(\vec{y},0)\right\} \,.
$$
As noticed by Olive and Witten \cite{Olive-Witten}, a non-zero
contribution to this comes from certain boundary terms which measure
electric and magnetic charges. To check this result, we have to look,
in the anti-commutator, for terms of the form $\int d^3 x\partial_i 
(A^{\dagger a} F^{aoi} + A^{\dagger a}\widetilde F^{aoi})$. For this,
we need only retain the relevant terms in the supercurrents:  
\begin{eqnarray}
g^2S^0_{(1)\alpha}&=&-i(\vec{\sigma}\sigma_y\lambda^{\dagger a}_1)_\alpha
\cdot (i\vec{F}^a+\vec{\widetilde F})^a +\sqrt{2}\,(\vec{\sigma}\cdot
\vec{D}A^{\dagger a}\lambda^a_{2})_\alpha + \cdots \,,\nonumber\\
g^2S^0_{(2)\alpha}&=&-i(\vec{\sigma}\sigma_y\lambda^{\dagger a}_2)_\alpha
\cdot (i\vec{F}^a+\vec{\widetilde F})^a -\sqrt{2}\,(\vec{\sigma}\cdot
\vec{D}A^{\dagger a}\lambda^a_{1})_\alpha + \cdots\,.
\end{eqnarray}
Using this, we get 
\begin{eqnarray}
\{Q_{(1)\alpha}\,,\,Q_{(2)\beta}\}&=&\frac{1}{g^4}\int d^3x\int d^3y\;
\big[i\sqrt{2}(\sigma^i\sigma_y)_{\alpha\gamma}\sigma^j_{\beta\lambda}
\{\lambda^{\dagger a}_{1\gamma},\lambda^b_{1\lambda}\}(iF^a_{0i}+
\widetilde F^a_{0i})D_j A^{\dagger b}     \nonumber\\
&&\qquad\qquad\quad -i\sqrt{2}(\sigma^j)_{\alpha\gamma}(\sigma^i
\sigma_y)_{\beta\lambda}\{\lambda^{\dagger a}_{2\gamma},
\lambda^b_{2\lambda}\}(iF^a_{0i}+\widetilde F^a_{0i})D_j 
A^{\dagger b}\big]      \nonumber \\
&=&\frac{1}{g^2}\int d^3x\,i\sqrt{2}\left[(\sigma^i\sigma_y
\sigma^{j\,T})_{\alpha\beta}-(\sigma^i\sigma_y\sigma^{j\,T})_{\beta\alpha}
\right](iF^a_{0i}+\widetilde F^a_{0i})D_j{A^\dagger}^{a}\,. \nonumber
\end{eqnarray}
The term within the square brackets, involving the $\sigma$ matrices,
can be simplified if we use $\sigma^i\sigma_y=-\sigma_y\sigma^{i\,T}$,
so that
$$
(\sigma^i\sigma_y\sigma^{j\,T})_{\alpha\beta}=-(\sigma_y\sigma^{i\,T} 
\sigma^{j\,T})_{\alpha\beta}=-[\sigma_y(\delta^{ij}-i\epsilon^{ijk}
\sigma_k^T)]_{\alpha\beta} \,.
$$
Subtracting from this a similar equation with $\alpha$ and $\beta$
interchanged, we get the term within the square brackets as equal to
$-2(\sigma_y)_{\alpha\beta}\delta^{ij}=2i\epsilon_{\alpha\beta}
\delta^{ij}$. Thus the commutator takes the form
$$
\{Q_{(1)\alpha},Q_{(2)\beta}\}=-\frac{2\sqrt{2}}{g^2}
\epsilon_{\alpha\beta}\int d^3x (iF^{a0i}+\widetilde F^{a0i}) 
D_i A^{\dagger a} \,.
$$
Using the Bianchi identity for the magnetic part, and the equation of
motion for the electric part, one can easily show that this is the
same as 
$$
\{Q_{(1)\alpha},Q_{(2)\beta}\}=-\frac{2\sqrt{2}}{g^2}
\epsilon_{\alpha \beta}\int d^3x\,\partial_i\left[(iF^{a0i}+
\widetilde F^{a0i}) A^{\dagger a}\right] \,.
$$
Similarly,
$$
\{\bar Q_{(1)\dot{\alpha}},\bar Q_{(2)\dot{\beta}}\}=-\frac
{2\sqrt{2}}{g^2} \epsilon_{\dot{\alpha}\dot{\beta}} \int d^3 x\,
\partial_i \left[(-iF^{a0i}+\widetilde F^{a0i}) A^a \right]\,.
$$
Thus, in the commutator of the supercurrents, we have recovered the
total derivatives which are nothing but the electric and magnetic
charge densities:
$$
Q_{ele}=-\frac{1}{ag}\int d^3x \partial_i(F^{a0i} A^a)= gn_e\,,\qquad 
Q_{mag}=-\frac{1}{ag}\int d^3x \partial_i(\widetilde F^{a0i} A^a)=
\frac{4\pi}{g}n_m \,.
$$
Here, $a$ is the value of $A$ in the Higgs vacuum. The charge
quantization condition used above corresponds to integral fundamantal
charges as in $SU(2)$ breaking to $U(1)$ with fileds in the adjoint
representation. Now, comparing with (\ref{Susy-C'}), we can easily
read off the $N=2$ central charge as $Z=-ia(n_e+ (4\pi
i/g^2)n_m)$. Recall that the phase of $Z$ is convention dependent and
we are mainly interested in its magnitude.  
 
To find the effect of the $\theta$-parameter, we can either modify
the calculation by adding the contribution from $\theta F\widetilde F$,
or simply use the Witten effect described in section (1.10). As we
learnt there, the effect of the $\theta$ parameter is to shift the
electric charge to $Q_{ele}=gn_e+(\theta g^2/8\pi^2)Q_{mag}$. Using
this (and ignoring an overall factor of $-i$), we get the central
charge as  
$$
Z= a\,(n_e + \tau_{cl} n_m)\,,\quad\rm{where},\quad 
\tau_{cl} = \frac{\theta}{2\pi} + \frac{4\pi i}{g^2}\,.
$$
Using supersymmetry algebra this implies a mass bound 
$$
M \geq \sqrt{2}\, |Z|= {\sqrt 2}\,|\,a\,(n_e + \tau_{cl} n_m )\,|\,,
$$
which is the BPS bound. So far, we have considered the microscopic
$N=2$ action (\ref{N=2}). At the level of the effective action
(\ref{N=2general}) specified by a prepotential ${\cal F}$, the central
charge takes the form 
$$
Z = an_e + a_D n_m \,,
$$
where $a_D=\partial{\cal F}/\partial a$. This formula can be
motivated as follows: As we will see in the next section, the theory
(\ref{N=2general}) has a dual description in which the magnetic
monopoles, and not the electric charges, appear as fundamental
objects. In this description, $n_e$ and $n_m$ are interchanged and
also $a$ is replaced by $a_D$. Thus, in the dual theory, one can
easily infer the contribution of the monopole to the BPS bound to be
$a_Dn_m$. The full duality group, as we will see in the next section,
is $SL(2,Z)$ under which $a$ and $a_D$ transform as a
doublet. Combining these facts leads to the central extension formula
given above. Similar constructions apply to arbitrary groups $SU(N)$
with adjoint matter.

\subsection{Central Charge in N=2 Gauge Theory with Matter}

When matter in the fundamental representation is added to the $N=2$
pure gauge theory described above, there is an important change: the
central charge also receives contributions from the masses of the
matter fields. In $N=2$ supersymmetry, matter fields are part of
hypermultiplets. A hypermultiplet can be most easily described in the
$N=1$ notation as consisting of a chiral superfield $Q$ and an
anti-chiral superfield $\widetilde Q^\dagger$. $Q$ ($\widetilde Q$)
contains components $q,\psi_q$ ($\widetilde q,\psi_{\widetilde q}$)
and transforms in the $N$ ($\bar{N}$) representation of the gauge
group $SU(N)$. Under the $SU(2)_{R}$ symmetry of the $N=2$ algebra,
$q,\widetilde q^\dagger$ form a doublet while $\psi_q$ and
$\psi^\dagger_{\tilde q}$ are singlets. Since the fields in a
hypermultiplet have spin $\leq 1/2$, they belong to a short
representation of the $N=2$ algebra and, therefore, must satify the
relation $M=\sqrt{2} Z$. The fact that $Z$, as derived in the previous
section, does not satisfy this relation in the presence of bare
hypermultiplet masses, indicates that the central charge must receive
further contributions from the hypermultiplets. We will calculate this
contribution below. 

Beside the standard kinetic terms and gauge couplings for the chiral
fields $Q$ and $\widetilde Q$, the $N=2$ Lagrangian in the presence of
matter also contains the $N=1$ superpotential given by 
$$
\sum_{i=1}^{N_f}\sqrt{2}\widetilde{Q}_i\Phi Q_i+\sum_{i=1}^{N_f} m_i 
\widetilde{Q}_i Q_i + h.c.
$$
Here, the index $i$ runs over the $N_f$ quark flavours and $\Phi$ is
the chiral superfield of the vector multiplet in the adjoint
representation. The first term is related to the gauge coupling of
the matter fields by $N=2$ and the second term is an $N=2$ invariant
mass term. When all masses are equal, the theory also has a $SU(N_f)$
flavour symmetry which is broken down to smaller subgroups when the
masses are not equal. For all masses unequal, it breaks down to
$U(1)^{N_f}$.  

The contribution of $Q$ and $\widetilde Q$ to the supercurrents can be
easily calculated using (\ref{N=1scch}), with the superpotential
${\cal W}$ as given above. Thus, as the contribution to
$S^\mu_{(1)\alpha}$, we get (supressing both colour and flavour
indices):
\begin{eqnarray}
\epsilon^\alpha S^\mu_{(1)\alpha}=\cdots&+&\sqrt{2}\,\epsilon\sigma^\nu  
\bar{\sigma}^\mu\psi_q D_\nu q^\dagger+i\sqrt{2}\,m\,
\widetilde{q}^\dagger \bar{\psi}_q\bar{\sigma}^\mu \epsilon\nonumber\\
&+& \sqrt{2}\,\epsilon\sigma^\nu\bar{\sigma}^\mu\psi_{\widetilde{q}}
D_\nu {\widetilde{q}^\dagger} + i \sqrt{2}\, m\, q^\dagger
\bar{\psi}_{\widetilde{q}} \bar{\sigma}^\mu \epsilon +\cdots\,.\nonumber
\end{eqnarray}
To obtain the corresponding terms in the second supercurrent, we have
to make the replacements $q\rightarrow\widetilde q^\dagger$, 
$\widetilde q^\dagger\rightarrow -q$,
\begin{eqnarray}
\epsilon^\alpha S^\mu_{(2)\alpha}=\cdots&+&\sqrt{2}\,\epsilon\sigma^\nu  
\bar{\sigma}^\mu\psi_q D_\nu \widetilde q - i\sqrt{2}\,m\,
q \bar{\psi}_q\bar{\sigma}^\mu \epsilon\nonumber\\
&-& \sqrt{2}\,\epsilon\sigma^\nu\bar{\sigma}^\mu\psi_{\widetilde{q}}
D_\nu q + i \sqrt{2}\, m\, \widetilde q \bar{\psi}_{\widetilde{q}} 
\bar{\sigma}^\mu \epsilon +\cdots\nonumber\,.
\end{eqnarray}
Using the standard canonical commutation relations and following a
procedure similar to the previous section, we can calculate the
contribution of these extra terms to the anti-commutator 
of the supercharges. The additional term is of the form
$2i\epsilon_{\alpha\beta}\,\sum_im_i S_i$, with
$$
S_i=\int d^3x\left(D_0q_i^\dagger q_i + q_iD_0q_i^\dagger
-\frac{i}{2}\psi^\dagger_{qi}\psi_{qi} + \frac{i}{2}\psi_{qi}
\psi^\dagger_{qi} -(q\rightarrow \widetilde q, \psi_q\rightarrow
\psi_{\widetilde q})\right)\,,
$$   
where, the index $i$ is not summed over. Clearly, this is a conserved
charge associated with a global $U(1)$ symmetry under which
$Q_i$ and $\widetilde Q_i$ carry charges $+1$ and $-1$,respectively. 
These are the $U(1)$ factors of the broken flavour group. Taking this
extra term into account, the formula for the central charge takes the
form 
\begin{equation}
Z = n_e a + n_m a_D + \sum_{i} \;\frac{1}{\sqrt{2}} \; m_i S_i \,.
\label{bps-m}
\end{equation}
This formula is crucial in the Seiberg-Witten analysis of $N = 2$
$SU(2)$ gauge theory with quark hypermultiplets.
